\begin{explanationbox}
\textbf{一、先看懂：焦张角究竟是什么？}

固定焦点 $F_1,F_2$，把椭圆上的点 $P$ 想成一名观察者：
从 $P$ 向两个焦点各拉一条线，眼前张开的角就是
$\theta=\angle F_1PF_2$。两焦点间距固定为 $2c$，
而椭圆定义又保证
\[
|PF_1|+|PF_2|=2a.
\]
所以我们研究的不是任意三角形，而是\textbf{底边固定、其余两边之和固定}的一族三角形。

\textbf{二、为什么短轴顶点看到的角最大？}

将 $P$ 沿椭圆上半部由 $(a,0)$ 移向 $(0,b)$：
开始时两条视线近乎重合，焦张角接近 $0$；随着 $P$ 远离焦点所在直线，
两条视线逐渐张开。到短轴顶点 $P=(0,b)$ 时，
$PF_1=PF_2=a$，焦点三角形恰好成为等腰三角形，张角最大。
之后 $P$ 移向 $(-a,0)$，过程反向。

\textbf{注意：} 这段是帮助观察的直觉，不以图形感觉代替证明。
真正保证角度随高度增大的精确关系是
\[
\boxed{\tan\frac\theta2=\frac{c\,h}{b^2}},
\qquad h=d(P,F_1F_2)=|y_P|.
\]
由于 $0<\theta/2<90^\circ$，$\tan(\theta/2)$ 随 $\theta$ 严格增大；
而椭圆上 $0\le h\le b$，所以 $h$ 最大的短轴顶点也使 $\theta$ 最大。

\textbf{三、把半角公式记成一句话}

\emph{半角的正切 $=$ 焦距参数 $c$ \(\times\) 到长轴的距离 $h$ \(/\) 短半轴平方 $b^2$。}
$P$ 上下翻转不会改变焦张角，因此必须用 $|y_P|$，而不是带符号的 $y_P$。
特别是 $P$ 越靠近长轴，$h$ 越小，$\tan(\theta/2)$ 就越小；
用这一点也能迅速判断倒写的公式不可能正确。

\textbf{四、为什么结果恰好是离心率 $e$？}

在短轴顶点，等腰三角形两腰长 $a$，底边一半长 $c$。
从 $P$ 向焦点连线作高，它同时平分顶角；
取其中一个直角三角形，直接得到
\[
\sin\frac{\theta_{\max}}2=\frac ca=e.
\]
这里的 $e$ 不是任意位置的 $\sin(\theta/2)$，
而是\textbf{半角正弦的最大值}。因此普通点只能得到
$\sin(\theta/2)\le e$。

\textbf{五、三种量之间如何互相转换？}
\begin{itemize}
  \item \textbf{知道离心率 $e$：} 立刻得到最大张角
  $\theta_{\max}=2\arcsin e$，并且椭圆上的非退化张角能取遍
  $(0,\theta_{\max}]$。
  \item \textbf{知道某个张角 $\theta_0$：} 只知道该角在椭圆上\emph{能出现}，
  就只能推出 $e\ge\sin(\theta_0/2)$；若题目说的是\emph{最大}张角，
  才能写等号。
  \item \textbf{知道点的高度 $h$：} 用
  $\tan(\theta/2)=ch/b^2$，不用分别求两条焦半径。
\end{itemize}

\textbf{六、与其他知识的连接}

本结论把椭圆定义（$r_1+r_2=2a$）、余弦定理、半角公式、
三角形面积 $S=c|y_P|$、三角形内切圆和均值不等式
$r_1r_2\le a^2$ 串成一条推理链。
思考同一题时，可以按\textbf{视角 $\rightarrow$ 三角形 $\rightarrow$ 周长/面积 $\rightarrow$ 半角}
这条路线识别突破口。
\end{explanationbox}
